\documentclass{book}

\title{Analisi Matematica T-A}
\author{Michele Mazza}
\date{November 2023}


% Pacchetti di Codifica e Lingua
\usepackage[utf8]{inputenc}
\usepackage[T1]{fontenc}
\usepackage[italian]{babel}

% Pacchetti Grafici e Diagrammi
\usepackage{graphicx}
\usepackage{tikz}
\usepackage{tikz-3dplot}
\usetikzlibrary{angles,quotes,arrows.meta, positioning}
\usepackage{pgfplots}
\pgfplotsset{compat=1.18}

% Pacchetti Matematici
\usepackage{amsmath, amssymb, amsthm}

% Pacchetti per Tabelle e Unità di Misura
\usepackage{geometry}
\usepackage{siunitx}
\usepackage{booktabs}
\usepackage{tabularx}
\usepackage{colortbl}

% Pacchetti per Layout e Formattazione
\usepackage{fancyhdr}
\usepackage{float}
\usepackage{microtype}

% Altri Pacchetti Utili
\usepackage{mdframed}
\usepackage{animate}

% Pacchetto enumitem
\usepackage{enumitem}

% Definizioni dei Colori
\definecolor{lightblue}{HTML}{9cd8fc}
\definecolor{myblue}{HTML}{139FE9}

% Definizione del comando \norm
\newcommand{\norm}[1]{\left\lVert#1\right\rVert}

%\font\nullfont=cmr10 

% Ambienti Personalizzati
\newmdtheoremenv[
    backgroundcolor=orange!10,
    linecolor=myblue,
    linewidth=0pt,
    leftmargin=0pt,
    rightmargin=0pt,
    innerleftmargin=10pt,
    innerrightmargin=10pt,
    innertopmargin=10pt,
    innerbottommargin=10pt,
    nobreak
]{curiosity}{Curiosità}

\newmdtheoremenv[
    font=\normalfont\small,
    linewidth=0pt,
    innerleftmargin=10pt,
    innerrightmargin=10pt,
]{example}{Esempio}

\newmdtheoremenv[
  backgroundcolor=lightblue,
  linecolor=myblue,
  linewidth=1.5pt,
  innertopmargin=10pt,
  innerbottommargin=10pt,
  innerrightmargin=10pt,
  innerleftmargin=10pt,
  nobreak
]{theorem}{Teorema}

\newmdtheoremenv[
  backgroundcolor=gray!10,
  linewidth=0pt,
  innertopmargin=10pt,
  innerbottommargin=10pt,
  innerrightmargin=10pt,
  innerleftmargin=10pt,
  nobreak
]{definition}{Definizione}

\newmdtheoremenv[
  backgroundcolor=gray!10,
  linecolor=myblue,
  linewidth=1.5pt,
  innertopmargin=10pt,
  innerbottommargin=10pt,
  innerrightmargin=10pt,
  innerleftmargin=10pt,
  nobreak
]{proposition}{Proposizione}

% Configurazioni di fancyhdr
\pagestyle{fancy}
\fancyhf{} % Pulisce tutti i campi
\fancyhead[LE,RO]{\thepage} % Numero di pagina
\fancyhead[RE]{\nouppercase{\leftmark}} % Titolo del capitolo
\fancyhead[LO]{\nouppercase{\rightmark}} % Titolo della sezione
\renewcommand{\headrulewidth}{0.5pt}
\renewcommand{\footrulewidth}{0pt}

\geometry{a4paper, left=20mm, right=20mm, top=20mm, bottom=20mm}

% Pacchetto Hyperref (inserito per ultimo)
\usepackage{hyperref}

\begin{document}

\frontmatter
\maketitle
\tableofcontents
\mainmatter

\chapter{Limiti e continuità}

\section{Intorni}
Nel definire i concetti di limite e di continuità, siamo condotti a considerare
numeri reali 'vicini' a un certo numero reale fissato, o, con linguaggio geometrico
equivalente, punti della retta 'vicini' a un punto fissato. Pertanto, iniziamo con
il precisare il concetto matematico di intorno di un punto.
\begin{definition}[\textbf{Intorno di un punto}]
    Sia $x_0 \in \mathbb{R}$ e sia $r > 0$. Si chiama \textbf{intorno} di $a$ di raggio $r$ l'insieme
    \begin{equation*}
        I_{r}(x_0) = \{x \in \mathbb{R} : |x - x_0| < r\}
    \end{equation*}
\end{definition}
\begin{figure}[H]
    \centering
    \begin{tikzpicture}
        \draw[-latex] (-1, 0) -- (5, 0);
        \filldraw [myblue] (2, 0) circle (2pt);
        \draw[myblue, thick] (0,0) -- (4, 0);
        \draw (0, -0.1) -- (0, 0.1);
        \draw (4, -0.1) -- (4, 0.1);
        \node [below] at (0, 0) {$x_0 - r$};
        \node [below] at (4, 0) {$x_0 + r$};
        \node [below] at (2, 0) {$x_0$};
        \node [above, myblue] at (2, 0.2) {$I_r$};
    \end{tikzpicture}
    \caption{Intorno di un punto}
    \label{fig:intorno}
\end{figure}
\begin{definition}[\textbf{Intorno di \(\pm\infty\)}]
    Sia $r > 0$. Si chiama \textbf{intorno} di $+\infty$ di \textbf{estremo inferiore} $r$ l'insieme
    \begin{equation*}
        I_{r}(+\infty) = \{x \in \mathbb{R} : x > r\}
    \end{equation*}
    e si chiama \textbf{intorno} di $-\infty$ di \textbf{estremo superiore} $-r$ l'insieme
    \begin{equation*}
        I_{r}(-\infty) = \{x \in \mathbb{R} : x < -r\}
    \end{equation*}
\end{definition}
\begin{figure}[H]
    \centering
    \begin{tikzpicture}
        \draw[-latex] (-1, 0) -- (5, 0);
        \filldraw (2, 0) circle (2pt);
        \draw[myblue, thick] (0,0) -- (-1, 0);
        \draw[myblue, thick] (4,0) -- (5, 0);
        \draw (0, -0.1) -- (0, 0.1);
        \draw (4, -0.1) -- (4, 0.1);
        \node [below] at (0, 0) {$-r$};
        \node [myblue, above] at (-0.5, 0.1) {$I_{r}(-\infty)$};
        \node [myblue, above] at (4.5, 0.1) {$I_{r}(+\infty)$};
        \node [below] at (-1, 0) {$-\infty$};
        \node [below] at (5, 0) {$+\infty$};
        \node [below] at (4, 0) {$r$};
        \node [below] at (2, 0) {$0$};
    \end{tikzpicture}
    \caption{Intorno di \(\pm\infty\)}
    \label{fig:intornoinfinito}
\end{figure}
\begin{definition}[\textbf{Punto di accumulazione di un insieme}]
    Sia $A \subseteq \mathbb{R}$ e sia $x_0 \in \mathbb{R}$. Si dice che $x_0$ è un \textbf{punto di accumulazione} per $A$ se ogni intorno di $x_0$ contiene almeno un punto di $A$ diverso da $x_0$. In simboli: \[\forall r > 0, I_r(x_0) \cap (A \setminus \{x_0\}) \neq \emptyset\]
\end{definition}
\begin{figure}[H]
    \centering
    \begin{tikzpicture}
        \draw[-latex] (-1, 0) -- (5, 0);
        \draw[ultra thick,darkgray] (0,0) -- (2, 0);
        \node [ultra thick, darkgray] at (0, 0) {[};
        \node [ultra thick, darkgray] at (2, 0) {)};
        \node [above, darkgray] at (1, 0.2) {$A$};
        \draw [myblue] (2, 0) circle (2pt);
        \node [above, myblue] at (2, 0.2) {$I_r$};
        \node [below, myblue] at (2, -0.1) {$x_0$};
        \node [ultra thick, myblue] at (1.5, 0) {(};
        \node [ultra thick, myblue] at (2.5, 0) {)};
    \end{tikzpicture}
    \caption{Punto di accumulazione}
    \label{fig:accumulazione}
\end{figure}
\begin{example}[\textbf{Punto di accumulazione for dummies}]
    Immagina una folla di persone che si radunano intorno a un famoso musicista in un parco. Il musicista è come un punto di accumulazione: non importa quanto ti avvicini, ci saranno sempre persone (punti) che si raccolgono intorno a lui, sempre più vicine. Ora, il bello è che il musicista potrebbe o non potrebbe essere parte di quella folla (l'insieme di persone): è possibile che sia al centro dell'attenzione, oppure che non sia neanche lì, magari è solo il posto dove si è esibito poco fa. In entrambi i casi, lui è il punto verso cui tutti si stanno avvicinando.
\end{example}
\noindent\begin{minipage}{\textwidth}
    \begin{definition}[\textbf{Derivato}]
        Si dice \textbf{derivato} di un un insieme \(A\), l'insieme dei punti di accumulazione di \(A\). In simboli:
        \[D(A) = \{x \in \mathbb{R} : x \text{ è un punto di accumulazione di } A\}\]
    \end{definition}
\end{minipage}
\begin{figure}[H]
    \centering
    \begin{tikzpicture}
        \filldraw [darkgray, opacity=0.1] (0,0) circle (2cm);
        \draw [dashed, myblue] (0,0) circle (2cm);
        \draw [dashed, myblue] (2,0) circle (0.5cm);
        \draw [myblue, opacity=0.1] (2,0) circle (0.5cm);
        \node [darkgray, above] at (0,2) {\textbf{A}=\{\( x,y \in \mathbb{R} : x^2+y^2<4\)\}};
        \node at (0,0) {\small\textbullet};
        \draw [myblue] (2,0)  circle (2pt);
        \node [right] at (0,0) {C};
        \node [myblue, right] at (2,0) {\(x_0\)};
        \node [myblue, right] at (2.5,0) {\(I_r\)};
        \node [myblue, below] at (0,-2) {\(D(A)=\{ x,y \in \mathbb{R} : x^2+y^2=4\}\)};
    \end{tikzpicture}
    \caption{Derivato di un insieme}
\end{figure}
\begin{theorem}[\textbf{Teorema di Bolzano - Weierstrass}]
    Sia \(A \subseteq \mathbb{R}\) finito e illimitato, allora \(A\) ha almeno un punto di accumulazione, ossia \(D(A) \neq \varnothing\)
\end{theorem}

\begin{definition}[\textbf{Punto isolato}]
    Sia $A \subseteq \mathbb{R}$ e sia $x_0 \in A$. Si dice che $x_0$ è un \textbf{punto isolato} per $A$ se esiste un intorno di $x_0$ che non contiene altri punti di $A$. In simboli: \[\exists r > 0 : I_r(x_0) \cap (A \setminus \{x_0\}) = \emptyset\]
\end{definition}
\begin{figure}[H]
    \centering
    \begin{tikzpicture}
        \draw[-latex] (-1, 0) -- (5, 0);
        \draw[ultra thick,darkgray] (0,0) -- (2, 0);
        \node [ultra thick, darkgray] at (0, 0) {[};
        \node [ultra thick, darkgray] at (2, 0) {)};
        \node [above, darkgray] at (1, 0.2) {$A$};
        \draw [myblue] (3, 0) circle (2pt);
        \node [above, myblue] at (3, 0.2) {$I_{r_1}$};
        \node [below, myblue] at (3, -0.1) {$x_0$};
        \node [ultra thick, myblue] at (2.5, 0) {(};
        \node [ultra thick, myblue] at (3.5, 0) {)};
        \node [ultra thick] at (1.5, 0) {(};
        \node [ultra thick] at (4.5, 0) {)};
        \node [above] at (3, 0.6) {$I_{r_2}$};
    \end{tikzpicture}
    \caption{Punto isolato}
    \begin{minipage}{0.8\textwidth}
        \small
        Anche se \(I_{r_2}\) soddisfa la definizione di punto di accumulazione, \(x_0\) non è un punto isolato perché \(I_{r_1}\) non contiene altri punti di \(A\) oltre a \(x_0\).
    \end{minipage}
    \label{fig:isolato}
\end{figure}

\begin{definition}[\textbf{Punto aderente}]
    Dato un insieme $A \subseteq \mathbb{R}$, si dice che $x_0 \in \mathbb{R}$ è un \textbf{punto aderente} per $A$ se ogni intorno di $x_0$ contiene almeno un punto di $A$. In simboli: \[\forall r > 0, I_r(x_0) \cap A \neq \emptyset\]
\end{definition}
\begin{figure}[H]
    \centering
    \begin{tikzpicture}
        \draw[-latex] (-1, 0) -- (5, 0);
        \draw[ultra thick,darkgray] (0,0) -- (2, 0);
        \node [ultra thick, darkgray] at (0, 0) {[};
        \node [ultra thick, darkgray] at (2, 0) {]};
        \node [above, darkgray] at (1, 0.2) {$A$};
        \filldraw [myblue] (2, 0) circle (1pt);
        \node [above, myblue] at (2, 0.2) {$I_r$};
        \node [below, myblue] at (2, -0.1) {$x_0$};
        \node [ultra thick, myblue] at (1.5, 0) {(};
        \node [ultra thick, myblue] at (2.5, 0) {)};
    \end{tikzpicture}
    \begin{tikzpicture}
        \draw[-latex] (-1, 0) -- (5, 0);
        \draw[ultra thick,darkgray] (0,0) -- (2, 0);
        \node [ultra thick, darkgray] at (0, 0) {[};
        \node [ultra thick, darkgray] at (2, 0) {]};
        \node [above, darkgray] at (1, 0.2) {$A$};
        \filldraw [myblue] (1.5, 0) circle (1pt);
        \node [above, myblue] at (1.5, 0.2) {$I_r$};
        \node [below, myblue] at (1.5, -0.1) {$x_0$};
        \node [ultra thick, myblue] at (1.1, 0) {(};
        \node [ultra thick, myblue] at (1.9, 0) {)};
    \end{tikzpicture}
    \caption{Punto aderente}
    \label{fig:aderente}
\end{figure}

\begin{curiosity}[Punto di accumulazione vs punto aderente]
    Vi è una differenza tra punto di accumulazione e punto aderente: il primo deve contenere almeno un punto di $A$ diverso da $x_0$, mentre il secondo può contenere anche $x_0$.
    \begin{figure}[H]
        \centering
        \begin{tikzpicture}
            \draw[-latex] (-1, 0) -- (5, 0);
            \draw[ultra thick,darkgray] (0,0) -- (2, 0);
            \node [ultra thick, darkgray] at (0, 0) {[};
            \node [ultra thick, darkgray] at (2, 0) {)};
            \node [above, darkgray] at (1, 0.2) {$A$};
            \draw [myblue] (2, 0) circle (2pt);
            \node [above, myblue] at (2, 0.2) {$I_r$};
            \node [below, myblue] at (2, -0.1) {$x_0$};
            \node [ultra thick, myblue] at (1.5, 0) {(};
            \node [ultra thick, myblue] at (2.5, 0) {)};
        \end{tikzpicture}
        \begin{tikzpicture}
            \draw[-latex] (-1, 0) -- (5, 0);
            \draw[ultra thick,darkgray] (0,0) -- (2, 0);
            \node [ultra thick, darkgray] at (0, 0) {[};
            \node [ultra thick, darkgray] at (2, 0) {]};
            \node [above, darkgray] at (1, 0.2) {$A$};
            \filldraw [myblue] (2, 0) circle (1pt);
            \node [above, myblue] at (2, 0.2) {$I_r$};
            \node [below, myblue] at (2, -0.1) {$x_0$};
            \node [ultra thick, myblue] at (1.5, 0) {(};
            \node [ultra thick, myblue] at (2.5, 0) {)};
        \end{tikzpicture}
        \caption{Punto di accumulazione vs punto aderente}
        \label{fig:accumulazionevsaderente}
    \end{figure}
\end{curiosity}

\begin{definition}[\textbf{Chiusura}]
    Si dice \textbf{chiusura} di un insieme $A \subseteq \mathbb{R}$ l'insieme dei punti aderenti ad $A$. In simboli: \[\overline{A} = \{x \in \mathbb{R} : x \text{ è un punto aderente di } A\} = A \cup D(A)\]
\end{definition}
\begin{figure}[H]
    \centering
    \begin{tikzpicture}
        \filldraw [myblue, opacity=0.1] (0,0) circle (2cm);
        \draw [myblue] (0,0) circle (2cm);
        \draw [dashed, myblue] (2,0) circle (0.5cm);
        \draw [dashed, myblue] (0,1) circle (0.5cm);
        \filldraw [myblue] (0,1)  circle (2pt);
        \node [myblue, right] at (0,1) {\(x_0\)};
        \node [myblue, right] at (0.5,1) {\(I_r\)};
        \draw [myblue, opacity=0.1] (2,0) circle (0.5cm);
        \node [darkgray, above] at (0,2) {\textbf{A}=\{\( x,y \in \mathbb{R} : x^2+y^2<4\)\}};
        \node at (0,0) {\small\textbullet};
        \filldraw [myblue] (2,0)  circle (2pt);
        \node [right] at (0,0) {C};
        \node [myblue, right] at (2,0) {\(x_0\)};
        \node [myblue, right] at (2.5,0) {\(I_r\)};
        \node [myblue, below] at (0,-2) {\(\overline{A}=\{ x,y \in \mathbb{R} : x^2+y^2\leq 4\}\)};
    \end{tikzpicture}
    \caption{Chiusura di un insieme}
    \label{fig:chiusura}
\end{figure}
\noindent\begin{minipage}{\textwidth}
    \begin{curiosity}[Derivato vs chiusura]
        Vi è una differenza tra derivato e chiusura: il primo è l'insieme dei punti di accumulazione, mentre il secondo è l'insieme dei punti aderenti.
        \begin{figure}[H]
            \centering
            \begin{tikzpicture}
                \filldraw [darkgray, opacity=0.1] (0,0) circle (2cm);
                \draw [dashed, darkgray] (0,0) circle (2cm);
                \draw [dashed, myblue] (0,0) circle (2cm);
                \draw [dashed, myblue] (2,0) circle (0.5cm);
                \draw [myblue, opacity=0.1] (2,0) circle (0.5cm);
                \node [darkgray, above] at (0,2) {\textbf{A}=\{\( x,y \in \mathbb{R} : x^2+y^2<4\)\}};
                \node at (0,0) {\small\textbullet};
                \draw [myblue] (2,0)  circle (2pt);
                \node [right] at (0,0) {C};
                \node [myblue, right] at (2,0) {\(x_0\)};
                \node [myblue, right] at (2.5,0) {\(I_r\)};
                \node [myblue, below] at (0,-2) {\(D(A)=\{ x,y \in \mathbb{R} : x^2+y^2=4\}\)};
            \end{tikzpicture}
            \begin{tikzpicture}
                \filldraw [myblue, opacity=0.1] (0,0) circle (2cm);
                \draw [myblue] (0,0) circle (2cm);
                \draw [dashed, myblue] (2,0) circle (0.5cm);
                \draw [dashed, myblue] (0,1) circle (0.5cm);
                \filldraw [myblue] (0,1)  circle (2pt);
                \node [myblue, right] at (0,1) {\(x_0\)};
                \node [myblue, right] at (0.5,1) {\(I_r\)};
                \draw [myblue, opacity=0.1] (2,0) circle (0.5cm);
                \node [darkgray, above] at (0,2) {\textbf{A}=\{\( x,y \in \mathbb{R} : x^2+y^2<4\)\}};
                \node at (0,0) {\small\textbullet};
                \filldraw [myblue] (2,0)  circle (2pt);
                \node [right] at (0,0) {C};
                \node [myblue, right] at (2,0) {\(x_0\)};
                \node [myblue, right] at (2.5,0) {\(I_r\)};
                \node [myblue, below] at (0,-2) {\(\overline{A}=\{ x,y \in \mathbb{R} : x^2+y^2\leq 4\}\)};
            \end{tikzpicture}
            \caption{Derivato vs chiusura}
            \label{fig:derivatovschiusura}
        \end{figure}
    \end{curiosity}
\end{minipage}
\noindent\begin{minipage}{\textwidth}
    \begin{definition}[\textbf{Punto di frontiera}]
        Sia $A \subseteq \mathbb{R}$ e sia $x_0 \in \mathbb{R}$. Si dice che $x_0$ è un \textbf{punto di frontiera} per $A$ se ogni intorno di $x_0$ contiene almeno un punto di $A$ e almeno un punto di $\mathbb{R} \setminus A$. In simboli: \[ \forall r > 0, I_r(x_0) \cap A \neq \emptyset \land I_r(x_0) \cap (\mathbb{R} \setminus A) \neq \emptyset\]
    \end{definition}
\end{minipage}
\begin{figure}[H]
    \centering
    \begin{tikzpicture}
        \draw[-latex] (-1, 0) -- (5, 0);
        \draw[ultra thick, darkgray] (0,0) -- (2, 0);
        \node [ultra thick, darkgray] at (0, 0) {[};
        \node [ultra thick, darkgray] at (2, 0) {]};
        \node [above, darkgray] at (1, 0.2) {$A$};
        \filldraw [myblue] (2, 0) circle (1pt);
        \filldraw [myblue] (0, 0) circle (1pt);
        \node [below, myblue] at (0, -0.1) {$x_0$};
        \node [above, myblue] at (0, 0.2) {$I_r$};
        \node [ultra thick, myblue] at (-0.5, 0) {(};
        \node [ultra thick, myblue] at (0.5, 0) {)};
        \node [above, myblue] at (2, 0.2) {$I_r$};
        \node [below, myblue] at (2, -0.1) {$x_0$};
        \node [ultra thick, myblue] at (1.5, 0) {(};
        \node [ultra thick, myblue] at (2.5, 0) {)};
    \end{tikzpicture}
    \caption{Punto di frontiera}
    \label{fig:frontiera}
\end{figure}
\begin{definition}[\textbf{Frontiera}]
    Si dice \textbf{frontiera} di un insieme $A \subseteq \mathbb{R}$ l'insieme dei punti di frontiera di $A$. In simboli: \[FrA = \partial A = \{x \in \mathbb{R} : x \text{ è un punto di frontiera di } A\}\]
\end{definition}
\begin{figure}[H]
    \centering
    \begin{tikzpicture}
        \filldraw [darkgray, opacity=0.1] (0,0) circle (2cm);
        \draw [myblue] (0,0) circle (2cm);
        \draw [dashed, myblue] (2,0) circle (0.5cm);
        \draw [myblue, opacity=0.1] (2,0) circle (0.5cm);
        \node [darkgray, above] at (0,2) {\textbf{A}=\{\( x,y \in \mathbb{R} : x^2+y^2<4\)\}};
        \node at (0,0) {\small\textbullet};
        \filldraw [myblue] (2,0)  circle (2pt);
        \node [right] at (0,0) {C};
        \node [myblue, right] at (2,0) {\(x_0\)};
        \node [myblue, right] at (2.5,0) {\(I_r\)};
        \node [myblue, below] at (0,-2) {\(FrA = \partial A =\{ x,y \in \mathbb{R} : x^2+y^2\ = 4\}\)};
    \end{tikzpicture}
    \caption{Frontiera di un insieme}
\end{figure}
\section{Limiti di successioni}
\noindent\begin{minipage}{\textwidth}
    \begin{definition}[\textbf{Successione}]
        Si dice \textbf{successione} una funzione $f : \mathbb{N} \rightarrow \mathbb{R}$, ossia una funzione che associa ad ogni numero naturale un numero reale. In simboli: \[f : \mathbb{N} \rightarrow \mathbb{R} \quad \text{oppure} \quad \{a_n\}_{n \in \mathbb{N}} \quad \text{oppure} \quad a : n \mapsto a_n\]
    \end{definition}
\end{minipage}
\begin{figure}[H]
    \centering
    \begin{tikzpicture}
        \begin{axis}[
                xlabel={\( x \)},
                ylabel={\( f(x) \)},
                xmin=0, xmax=5,
                ymin=0, ymax=3,
                grid=both,
                axis lines=middle,
                axis line style=->,
                grid style={line width=.1pt, draw=gray!10},
                major grid style={line width=.2pt,draw=gray!50},
                minor tick num=5,
                enlargelimits={abs=0.5},
                ticklabel style={font=\small},
                xlabel style={at={(ticklabel* cs:1)}, anchor=north west},
                ylabel style={at={(ticklabel* cs:1)}, anchor=south west},
                domain=0:5,
                samples=1000,
                unbounded coords=jump, % Gestisce i punti non definiti
                legend pos=south east
            ]
            \addplot [myblue,only marks] coordinates {(0,1) (1,2) (2,2.25) (3,2.37) (4,2.44) (5,2.49)};
            \addplot [dashed, domain=0:6] {3};
            \addplot [dashed, domain=0:6] {2};
            \addplot [dashed,myblue, domain=0:6] {e};
            \node [left, myblue] at (axis cs: 0,e) {\( e \)};
            \legend{\(a_n=(1+\frac{1}{n})^n\)}
        \end{axis}
    \end{tikzpicture}
    \caption{Successione}
    \label{fig:successione}
\end{figure}
\begin{definition}[\textbf{Limite convergente di una successione}]
    Si dice che la successione \( a : n \mapsto a_n \) \textbf{tende al limite} \(l \in \mathbb{R}\) (oppure \textbf{converge} a \(l\), oppure \textbf{ha limite} \(l\)), e si scrive
    \[\lim_{n \to +\infty} a_n = l\]
    se
    \[
        \begin{cases}
            \forall \varepsilon > 0, \text{esiste un intero } n_\varepsilon \text{ tale che} \\
            \forall n \geq n_0, \quad n>n_\varepsilon \implies |a_n - l| < \varepsilon       \\
        \end{cases}
    \]
    in simboli:
    \[
        \forall \varepsilon > 0, \exists n_\varepsilon \in \mathbb{N} : \forall n \geq n_\varepsilon, n > n_\varepsilon \implies |a_n - l| < \varepsilon
    \]
\end{definition}
\begin{curiosity}[Estensione a \(\pm\infty\)]
    Con la terminologia degli intorni la condizione precedente si può scrivere come \(n \in I_\varepsilon(+\infty)\), mentre \(|a_n - l| < \varepsilon\) equivale a \(a_n \in I_\varepsilon(l)\).
    Pertanto, la definizione di limite può essere riscritta come
    \[
        \forall I_\varepsilon(l), \exists I_{n_\varepsilon}(+\infty) : \forall n \geq n_\varepsilon, n \in I_\varepsilon(+\infty) \implies a_n \in I_\varepsilon(l)
    \]
\end{curiosity}
\begin{definition}[\textbf{Limite divergente di una successione}]
    Si dice che la successione \( a : n \mapsto a_n \) \textbf{tende al limite} \(+\infty\) (oppure \textbf{diverge a \(+\infty\)}), e si scrive
    \[\lim_{n \to +\infty} a_n = +\infty\]
    se
    \[
        \begin{cases}
            \forall \text{ numero reale } A > 0, \text{esiste un intero } n_A \text{ tale che} \\
            \forall n \geq n_0, \quad n>n_A \implies a_n > A                                   \\
        \end{cases}
    \]
    in simboli:
    \[
        \forall I_A(l), \exists I_{n_A}(+\infty) : \forall n \geq n_0, n \in I_A(+\infty) \implies a_n \in I_A(l)
    \]
\end{definition}
\begin{proposition}
    Una successione può essere \textbf{convergente} o \textbf{divergente} a \(\pm\infty\), se non è né convergente né divergente allora è \textbf{indeterminata} o \textbf{oscillante}.
\end{proposition}

\begin{theorem}[\textbf{Condizione sufficiente per escludere il comportamento indeterminato di una successione è la monotonia}]
    Sia \(a_n\) una successione monotona. Allora essa è convergente oppure divergente a \(\pm\infty\). Precisamente, nel caso in cui la successione sia crescente si ha:
    \begin{itemize}
        \item se \(a_n\) è limitata superiormente, allora converge al \(\sup a_n \)
        \item se \(a_n\) non è limitata superiormente, allora diverge a \(+\infty\)
    \end{itemize}
    Se invece la successione è decrescente la definizione si modifica in modo ovvio.
\end{theorem}
\subsection{\textbf{Il numero di Nepero o Eulero e}}
\noindent\begin{minipage}{\textwidth}
    \begin{theorem}[\(a_n\) è crescente e limitata superiormente]
        Sia \(a_n=(1+\frac{1}{n})^n\). Allora:
        \begin{enumerate}
            \item \(a_n\) è crescente
            \item \(a_n\) è limitata superiormente
        \end{enumerate}
        \textbf{Dimostrazione}
        Per verificare che la successione \(a_n=(1+\frac{1}{n})^n\) sia crescente e limitata bisogna prima introdurre:
        \begin{itemize}
            \item  il \textbf{binomio di Newton}:
                  \begin{equation*}
                      (a+b)^n = \sum_{k=0}^{n} \binom{n}{k} a^{n-k}b^k
                  \end{equation*}
            \item il \textbf{coefficiente binomiale}:
                  \begin{equation*}
                      \binom{n}{k} = \frac{n!}{k!(n-k)!}
                  \end{equation*}
            \item la \textbf{Disuguaglianza di Bernoulli}:
                  \begin{equation*}
                      (1+x)^n \geq 1+nx \quad \forall x \geq -1
                  \end{equation*}
        \end{itemize}
        \begin{enumerate}
            \item Per dimostrare che \(a_n\) è monotona strettamente crescente, dobbiamo mostrare che \(a_n < a_{n+1}\) per ogni \(n\), è però più conveniente dimostrare che \(a_{n-1} < a_n\) che è la stessa cosa.
            \item Per poi dimostrare che \(a_n\) è limitata superiormente, dobbiamo mostrare che \(a_n < 3\) per ogni \(n\).
        \end{enumerate}
    \end{theorem}
\end{minipage}

\section{Limiti di funzioni}

\subsection{Limiti finiti}

\subsubsection{\texorpdfstring{Caso \(x_0 \in \mathbb{R}\), \(L \in \mathbb{R}\)}{Caso x_0 in R, L in R}}

\begin{definition}[Caso \(x_0 \in \mathbb{R}\), \(l \in \mathbb{R}\)]
    \[
        \lim_{x \to x_0} f(x) = l
    \]
    se
    \[
        \forall \varepsilon > 0, \, \exists \delta > 0 \ : \ |f(x) - l| < \varepsilon, \forall x \in (x_0 - \delta, x_0 + \delta) \setminus \{x_0\}
    \]
    \textbf{Osservazioni}
    \begin{itemize}
        \item \(x_0\) è un punto di accumulazione per \(domf\), dunque \(x_0\) non è necessariamente in \(domf\).
    \end{itemize}
\end{definition}

\begin{figure}[H]
    \centering
    \begin{tikzpicture}
        % grid and Axis
        \draw[help lines, color=gray!30, dashed] (-3 , -3) grid (5 , 4);
        \draw[->] (-3 , 0) -- (5 , 0) node[right]{$x$};
        \draw[->] (0 , -3) -- (0, 4) node[above]{$y$};
        % Bounds and function value
        \draw[dashed] (3.5, 0) node[below right] {$\scriptstyle x_0+\delta$} -- (3.5 , 3.1) -- (0 , 3.1) node[left] {$\scriptstyle l+\varepsilon$}; %Upper
        \draw[dashed,] (3  , 0) node[below = 1.5 mm] {$\scriptstyle x_0$} -- (3 , 2.55) -- (0 , 2.55) node[left] {$\scriptstyle l$}; % l
        \draw[dashed] (2.5 ,  0) node[below left] {$\scriptstyle x_0-\delta$} -- (2.5 , 2.0) -- (0 , 2.0) node[left] {$\scriptstyle l-\varepsilon$}; % Lower
        \draw[thick,myblue,domain=-1.8 : 2.25 , samples = 20] {} plot ((\x + 2 , {exp(\x + 1.5)/(2 +  exp(\x)}); % f(x)
        % Annotations
        \draw [|-|, thick, myblue ] (2.5 , 0) -- (3.5 , 0) node[midway, below = 3mm]{$\scriptstyle I(x_0)$};
        \draw [<-> , thin] (0.25 , 2.55) -- (0.25 , 3.1) node[midway, right] {$\scriptstyle \varepsilon$}; % Upper Arrow
        \draw [<-> , thin] (0.25 , 2) -- (0.25 , 2.55)  node[midway, right] {$\scriptstyle \varepsilon$}; % Lower Arrow    
        \filldraw [white] (3 , 2.55) circle (1.5pt);
        \draw [black] (3 , 2.55) circle (1.5pt) node[above right] {$\scriptstyle (x_0,l)$};
    \end{tikzpicture}
    \caption{Limite finito per \(x\) che tende a \(x_0\)}
    \label{fig:limitefinito}
\end{figure}

\subsubsection{\texorpdfstring{Caso \(x_0 = \pm\infty\), \(L \in \mathbb{R}\)}{Caso x_0 = pm infinito, L in R}}

\begin{definition}[Caso \(x_0 = \pm\infty\), \(l \in \mathbb{R}\)]
    \[
        \lim_{x \to \pm\infty} f(x) = l
    \]
    se
    \[
        \forall \varepsilon > 0, \, \exists M > 0 \ : \ |f(x) - l| < \varepsilon, \forall x > M \ (\text{per } x_0 = +\infty) \ \text{o} \ \forall x < -M \ (\text{per } x_0 = -\infty)
    \]
\end{definition}
\begin{figure}[H]
    \centering
    \begin{tikzpicture}[
            declare function={
                    F(\x) = 1/(\x);
                }]
        \begin{axis}
            [
                xlabel={\( x \)},
                ylabel={\( f(x) \)},
                xmin=-5, xmax=5,
                ymin=-5, ymax=5,
                grid=both,
                axis lines=middle,
                axis line style=->,
                grid style={line width=.1pt, draw=gray!10},
                major grid style={line width=.2pt,draw=gray!50},
                minor tick num=5,
                enlargelimits={abs=0.5},
                xlabel style={at={(ticklabel* cs:1)}, anchor=north west},
                ylabel style={at={(ticklabel* cs:1)}, anchor=south west},
                samples=100,
                unbounded coords=jump, % Gestisce i punti non definiti
                legend pos=south east
            ]
            \addplot[very thick,myblue,domain=-5:-0.01]{F(x)};
            \addplot[very thick,myblue,domain=0.01:5]{F(x)};
            \legend{{\(f(x)=\frac{1}{x}\)}}
        \end{axis}
    \end{tikzpicture}
    \caption{Limite finito per \(x\) che tende a \(+\infty\)}
    \label{fig:limitefinitoinfinito}
\end{figure}
In questo caso \ref{fig:limitefinitoinfinito} \(\lim_{x \to +\infty} f(x) = 0 \).

\noindent\begin{minipage}{\textwidth}
    \begin{curiosity}
        In generale, il fatto che una funzione abbia limite finito in un punto \(x_0\) non implica che la funzione sia definita in \(x_0\).
        Sono possibili i seguenti casi:
        \begin{itemize}
            \item \(\exists \quad \lim_{x \to x_0} f(x) = l \quad e \quad l = f(x_0)\)
            \item \(\exists \quad \lim_{x \to x_0} f(x) = l \quad e \quad l \neq f(x_0)\)
            \item \(\exists \quad \lim_{x \to x_0} f(x) = l \quad e \quad \nexists \quad f(x_0)\)
        \end{itemize}
    \end{curiosity}
\end{minipage}
\subsection{Limiti infiniti}

\subsubsection{\texorpdfstring{Caso \(x_0 \in \mathbb{R}\), \(L = \pm \infty\)}{Caso x_0 in R, L = pm infinito}}

\begin{definition}[Caso \(x_0 \in \mathbb{R} \), \(L = \pm \infty\)]
    \[
        \lim_{x \to x_0} f(x) = L
    \]
    se
    \[
        \forall A > 0, \, \exists \delta > 0 \ : \ (\text{se } L = +\infty, \ f(x) > A \ \text{o se } L = -\infty, \ f(x) < -A), \forall x \in (x_0 - \delta, x_0 + \delta) \setminus \{x_0\}
    \]
\end{definition}
\begin{figure}[H]
    \centering
    \begin{tikzpicture}[
            declare function={
                    F(\x) = 1/(\x);
                }]
        \begin{axis}
            [
                xlabel={\( x \)},
                ylabel={\( f(x) \)},
                xmin=-5, xmax=5,
                ymin=-5, ymax=5,
                grid=both,
                axis lines=middle,
                axis line style=->,
                grid style={line width=.1pt, draw=gray!10},
                major grid style={line width=.2pt,draw=gray!50},
                minor tick num=5,
                enlargelimits={abs=0.5},
                xlabel style={at={(ticklabel* cs:1)}, anchor=north west},
                ylabel style={at={(ticklabel* cs:1)}, anchor=south west},
                samples=100,
                unbounded coords=jump, % Gestisce i punti non definiti
                legend pos=south east
            ]
            \addplot[very thick,myblue,domain=-5:-0.01]{-F(x)};
            \addplot[very thick,myblue,domain=0.01:5]{F(x)};
            \legend{{\(f(x)=\left|\frac{1}{x}\right|\)}}
        \end{axis}
    \end{tikzpicture}
    \caption{Limite infinito per \(x\) che tende a \(x_0\)}
    \label{fig:limiteinfinito}
\end{figure}
In questo caso \ref{fig:limiteinfinito} \(\lim_{x \to 0} f(x) = +\infty\), ma \(f(0)\) non è definita.

\subsubsection{\texorpdfstring{Caso \(x_0 = \pm \infty \), \(L = \pm \infty\)}{Caso x_0 = pm infinito, L = pm infinito}}

\begin{definition}[Caso \(x_0 = \pm \infty \), \(L = \pm \infty\)]
    \[
        \lim_{x \to \pm \infty} f(x) = L
    \]
    se
    \[
        \forall A > 0, \, \exists M > 0 \ : \ (\text{se } L = +\infty, \ f(x) > A \ \text{o se } L = -\infty, \ f(x) < -A), \forall x \ : \ |x| > M
    \]
\end{definition}
\begin{figure}[H]
    \centering
    \begin{tikzpicture}[
            declare function={
                    F(\x) = \x^3;
                }]
        \begin{axis}
            [
                xlabel={\( x \)},
                ylabel={\( f(x) \)},
                xmin=-5, xmax=5,
                ymin=-5, ymax=5,
                grid=both,
                axis lines=middle,
                axis line style=->,
                grid style={line width=.1pt, draw=gray!10},
                major grid style={line width=.2pt,draw=gray!50},
                minor tick num=5,
                enlargelimits={abs=0.5},
                xlabel style={at={(ticklabel* cs:1)}, anchor=north west},
                ylabel style={at={(ticklabel* cs:1)}, anchor=south west},
                samples=100,
                unbounded coords=jump, % Gestisce i punti non definiti
                legend pos=south east
            ]
            \addplot[very thick,myblue,domain=-5:-0.01]{F(x)};
            \addplot[very thick,myblue,domain=0.01:5]{F(x)};
            \legend{{\(f(x)=x^3\)}}
        \end{axis}
    \end{tikzpicture}
    \caption{Limite infinito per \(x\) che tende a \(\pm\infty\)}
    \label{fig:limiteinfinitoinfinito}
\end{figure}
In questo caso \ref{fig:limiteinfinitoinfinito} \(\lim_{x \to +\infty} f(x) = +\infty\).

\section{Continuità di una funzione}
\begin{definition}[\textbf{Funzione continua in \(x_0\)}]
    Una f(x) è \textbf{continua} in \(x_0\) se
    \[\lim_{x \to x_0} f(x) = f(x_0)\]
    ovvero
    \[\lim_{x \to x_{0^-}} f(x) = \lim_{x \to x_{0^+}} f(x) = f(x_0)\]
\end{definition}
\begin{definition}
    [\textbf{Funzione continua in un intervallo}]
    Una funzione \(f(x)\) è \textbf{continua} in un intervallo \([a,b]\) se è continua in ogni punto di \([a,b]\).
\end{definition}
\begin{proposition}
    Tutte le \textbf{funzioni elementari} (polinomi e funzioni razionali,
    funzioni elevamento a potenza, funzioni trigonometriche, funzioni
    esponenziali e le loro funzioni inverse) \textbf{sono continue} in tutto il loro
    dominio.
\end{proposition}
\subsection{Tipologie di discontinuità}
\begin{definition}[\textbf{Discontinuità di prima specie}]
    Sia \(f(x)\) definita in un intorno di \(x_0\). Allora \(f(x)\) è \textbf{discontinua di prima specie} in \(x_0\) se
    \[\lim_{x \to x_{0^-}} f(x) \neq \lim_{x \to x_{0^+}} f(x)\]
    Tale discontinuità è detta anche \textbf{a salto}, la lunghezza di tale salto è:
    \[
        s = |l-m|
    \]
\end{definition}
\begin{figure}[H]
    \centering
    \begin{tikzpicture}[
            declare function={
                    sigma(\x)=1/(1+exp(-\x));
                    F(\x) = 0.5 * sigma(\x);
                }]
        \begin{axis}
            [
                xlabel={\( x \)},
                ylabel={\( f(x) \)},
                xmin=0, xmax=2,
                ymin=0, ymax=2,
                grid=both,
                axis lines=middle,
                axis line style=->,
                grid style={line width=.1pt, draw=gray!10},
                major grid style={line width=.2pt,draw=gray!50},
                minor tick num=5,
                enlargelimits={abs=0.5},
                xlabel style={at={(ticklabel* cs:1)}, anchor=north west},
                ylabel style={at={(ticklabel* cs:1)}, anchor=south west},
                samples=100,
                unbounded coords=jump, % Gestisce i punti non definiti
                legend pos=south east
            ]
            \draw [gray, dashed] (1,0.36) -- (0,0.36) node [left] {\(l\)};
            \draw [gray, dashed] (1,0.86) -- (0,0.86) node [left] {\(m\)};
            \addplot[very thick,myblue,domain=-1:1]   {F(x)};
            \addplot[very thick,myblue,domain=1:3]    {F(x)+0.5};
            \draw [dashed,black!50] (1,{F(1)}) -- (1,{F(1)+0.5});
            \filldraw [myblue] (1,0.36) circle (2pt);
            \filldraw [white] (1,0.86) circle (2pt);
            \draw [myblue] (1,0.86) circle (2pt);
            \node [myblue, below] at (1,0.36) {\(x_0\)};

        \end{axis}
    \end{tikzpicture}
    \caption{Discontinuità di prima specie}
    \label{fig:discontinuitaprimaspecie}
\end{figure}
\noindent\begin{minipage}{\textwidth}
    \begin{definition}
        [\textbf{Discontinuità di seconda specie}]
        Sia \(f(x)\) definita in un intorno di \(x_0\). Allora \(f(x)\) è \textbf{discontinua di seconda specie} in \(x_0\) se
        \[\lim_{x \to x_{0^-}} f(x) = \pm \infty\ \lor\ \nexists \quad \lor \quad \lim_{x \to x_{0^+}} f(x) = \pm \infty\ \lor\ \nexists\]
    \end{definition}
\end{minipage}
\begin{figure}[H]
    \centering
    \begin{tikzpicture}[
            declare function={
                    F(\x) = 1/(\x);
                }]
        \begin{axis}
            [
                xlabel={\( x \)},
                ylabel={\( f(x) \)},
                xmin=-5, xmax=5,
                ymin=-5, ymax=5,
                grid=both,
                axis lines=middle,
                axis line style=->,
                grid style={line width=.1pt, draw=gray!10},
                major grid style={line width=.2pt,draw=gray!50},
                minor tick num=5,
                enlargelimits={abs=0.5},
                xlabel style={at={(ticklabel* cs:1)}, anchor=north west},
                ylabel style={at={(ticklabel* cs:1)}, anchor=south west},
                samples=100,
                unbounded coords=jump, % Gestisce i punti non definiti
                legend pos=south east
            ]
            \addplot[very thick,myblue,domain=-5:-0.01]{F(x)};
            \addplot[very thick,myblue,domain=0.01:5]{F(x)};
            \filldraw [myblue] (0,0) circle (2pt) node[below left] {\(x_0\)};
            \node[draw, align=center, fill=white, scale= 0.8] (legend) at (3.1,-3.9) {
                \(f(x) = \begin{cases}
                    \frac{1}{x} & \text{se } x \neq 0 \\
                    0           & \text{se } x = 0
                \end{cases}\)
            };
        \end{axis}
    \end{tikzpicture}
    \caption{Discontinuità di seconda specie}
    \label{fig:discontinuitasecondaspecie}
\end{figure}
\begin{definition}
    [\textbf{Discontinuità di terza specie}]
    Sia \(f(x)\) definita in un intorno di \(x_0\). Allora \(f(x)\) è \textbf{discontinua di terza specie} in \(x_0\) se
    \[\lim_{x \to x_{0}} f(x) = l \quad f(x_0) \neq l\]
    Tale discontinuità è anche detta \textbf{eliminabile}.
\end{definition}
\begin{figure}[H]
    \centering
    \begin{tikzpicture}[
            declare function={
                    F(\x) = x;
                }]
        \begin{axis}
            [
                xlabel={\( x \)},
                ylabel={\( f(x) \)},
                xmin=-2, xmax=2,
                ymin=-2, ymax=2,
                grid=both,
                axis lines=middle,
                axis line style=->,
                grid style={line width=.1pt, draw=gray!10},
                major grid style={line width=.2pt,draw=gray!50},
                minor tick num=5,
                enlargelimits={abs=0.5},
                xlabel style={at={(ticklabel* cs:1)}, anchor=north west},
                ylabel style={at={(ticklabel* cs:1)}, anchor=south west},
                samples=100,
                unbounded coords=jump, % Gestisce i punti non definiti
                legend pos=south east
            ]
            \addplot[very thick,myblue,domain=-2:2]{F(x)};
            \draw [gray, dashed] (1,1) -- (1,0);
            \filldraw [myblue] (1,0) circle (2pt) node[above right] {\(x_0\)};
            \filldraw [white] (1,1) circle (2pt);
            \draw [myblue] (1,1) circle (2pt);
            \node[draw, align=center, fill=white, scale= 0.8] (legend) at (1.4,-1.9) {
                \(f(x) = \begin{cases}
                    x & \text{se } x \neq 1 \\
                    0 & \text{se } x = 1
                \end{cases}\)
            };
        \end{axis}
    \end{tikzpicture}
    \caption{Discontinuità di terza specie}
    \label{fig:discontinuitaterzaspecie}
\end{figure}

\subsection{Teoremi sulle funzioni continue}

\begin{theorem}[\textbf{Teorema dei valori intermedi}]
    Sia \(f(x)\) continua in \([a,b]\). Allora \(f(x)\) assume tutti i valori compresi tra \(f(a)\) e \(f(b)\).
\end{theorem}

\begin{theorem}[\textbf{Teorema della composizione di funzioni continue}]
    Siano \(f(x)\) continua in \(x_0\) e \(g(x)\) continua in \(f(x_0)\). Allora \(g \circ f\) è continua in \(x_0\).
\end{theorem}

\begin{theorem}[\textbf{Teorema della somma di funzioni continue}]
    Siano \(f(x)\) e \(g(x)\) continue in \(x_0\). Allora \(f(x) + g(x)\) è continua in \(x_0\).
\end{theorem}

\begin{theorem}[\textbf{Teorema del prodotto di funzioni continue}]
    Siano \(f(x)\) e \(g(x)\) continue in \(x_0\). Allora \(f(x) \cdot g(x)\) è continua in \(x_0\).
\end{theorem}

\begin{theorem}[\textbf{Teorema di Weierstrass}]
    Sia \(f(x)\) continua in \([a,b]\). Allora \(f(x)\) è limitata e ammette massimo e minimo in \([a,b]\).
\end{theorem}

\begin{theorem}[\textbf{Teorema di Bolzano-Weierstrass}]
    Sia \(f(x)\) continua in \([a,b]\). Allora \(f(x)\) assume il massimo e il minimo in \([a,b]\).
\end{theorem}
\end{document}